\documentclass[aspectratio=169]{beamer}
\usepackage[utf8]{inputenc}
\usepackage{booktabs}
\usepackage{listings}
\usepackage{xcolor}
\usepackage{colortbl}
\usepackage{graphicx}
\usepackage{amsmath}

% Theme Setup to match CMU-DB style
\usetheme{Madrid}
\usecolortheme{beaver} % Red/Gray theme roughly matches
\setbeamercolor{title}{bg=darkgray,fg=white}
\setbeamercolor{frametitle}{bg=darkgray,fg=white}
\setbeamercolor{structure}{fg=darkred}

% SQL Syntax Highlighting
\definecolor{codegreen}{rgb}{0,0.6,0}
\definecolor{codegray}{rgb}{0.5,0.5,0.5}
\definecolor{codepurple}{rgb}{0.58,0,0.82}
\definecolor{backcolour}{rgb}{0.95,0.95,0.92}

\lstdefinestyle{sqlstyle}{
    backgroundcolor=\color{backcolour},   
    commentstyle=\color{codegreen},
    keywordstyle=\color{blue}\bfseries,
    numberstyle=\tiny\color{codegray},
    stringstyle=\color{codepurple},
    basicstyle=\ttfamily\footnotesize,
    breakatwhitespace=false,         
    breaklines=true,                 
    captionpos=b,                    
    keepspaces=true,                 
    numbers=none,                    
    numbersep=5pt,                  
    showspaces=false,                
    showstringspaces=false,
    showtabs=false,                  
    tabsize=2,
    upquote=true,
    language=SQL
}
\lstset{style=sqlstyle}

% Meta Data
\title{Modern SQL}
\subtitle{15-445/645 -- Fall 2024}
\author{Prof. Andy Pavlo}
\institute{Carnegie Mellon University}
\date{Fall 2024}

\begin{document}

%--- Slide 1 ---
\frame{\titlepage}

%--- Slide 2 ---
\begin{frame}{Today's Agenda}
    \begin{enumerate}
        \item Database Systems Background
        \item Relational Model
        \item Relational Algebra
        \item Alternative Data Models
        \item Q\&A Session
    \end{enumerate}
\end{frame}

%--- Slide 3 ---
\begin{frame}{Last Class}
    \begin{itemize}
        \item We introduced the \textbf{Relational Model} as the superior data model for databases.
        \item We then showed how \textbf{Relational Algebra} is the building blocks that will allow us to query and modify a relational database.
    \end{itemize}
\end{frame}

%--- Slide 4-5 ---
\begin{frame}{SQL History}
    \begin{itemize}
        \item In 1971, IBM created its first relational query language called \textbf{SQUARE}.
        \item IBM then created "\textbf{SEQUEL}" in 1972 for IBM System R prototype DBMS.
        \begin{itemize}
            \item Structured English Query Language
        \end{itemize}
        \item IBM releases commercial SQL-based DBMSs:
        \begin{itemize}
            \item System/38 (1979)
            \item SQL/DS (1981)
            \item DB2 (1983)
        \end{itemize}
    \end{itemize}
    \vspace{0.5cm}
    \textbf{Original Syntax Examples (System R Papers):}
    \begin{itemize}
        \item Q2. Find average salary of employees in Shoe Dept:
        \[ AVG \: ({SAL}^{EMP'}_{DEPT}('SHOE')) \]
        \item Q3. Find items sold by departments on the 2nd floor:
        \[ SALES \cdot ITEM \: (LOC \cdot DEPT \: (DEPT \cdot FLOOR ('2'))) \]
    \end{itemize}
\end{frame}

%--- Slide 6 ---
\begin{frame}{SQL Standards History}
    \begin{itemize}
        \item \textbf{ANSI Standard} in 1986. \textbf{ISO} in 1987.
        \item Current standard is \textbf{SQL:2023}.
    \end{itemize}
    \begin{description}
        \item[SQL:2023] Property Graph Queries, Multi-Dim. Arrays
        \item[SQL:2016] JSON, Polymorphic tables
        \item[SQL:2011] Temporal DBs, Pipelined DML
        \item[SQL:2008] Truncation, Fancy Sorting
        \item[SQL:2003] XML, Windows, Sequences, Auto-Gen IDs
        \item[SQL:1999] Regex, Triggers, OO
    \end{description}
    \vspace{0.2cm}
    \textit{The minimum language syntax a system needs to say that it supports SQL is \textbf{SQL-92}.}
\end{frame}

%--- Slide 7-8 ---
\begin{frame}{SQL History: The "Death of SQL" Predictions}
    \begin{block}{Jo Kristian Bergum (@jobergum) - Apr 27, 2023}
    "Tensor and vector databases will replace most legacy databases in this decade. A disruption fueled by natural language interfaces and deep neural representations. In other words: \textbf{Natural query languages (NQL) replace the structured query language (SQL).}"
    \end{block}
    
    \begin{block}{Gagan Biyani (@gaganbiyani) - May 18, 2023}
    "SQL is going to die at the hands of an AI. I'm serious. @mayowaoshin is already doing this. Takes your company's data and ingests it into ChatGPT... just ask it questions using natural language."
    \end{block}
\end{frame}

%--- Slide 9 ---
\begin{frame}{SQL History: The Reality (IEEE Spectrum)}
    \begin{columns}
        \begin{column}{0.6\textwidth}
           \textbf{IEEE Spectrum: The Rise of SQL} (Aug 2022)
           \begin{itemize}
               \item "It's become the second programming language everyone needs to know."
               \item SQL has become the primary query language for accessing and managing data specifically in relational databases.
               \item \textbf{Rankings:} SQL dominated the jobs ranking in IEEE Spectrum's interactive rankings.
               \item "Normally top position is occupied by Python... but the sheer number of times employers said they wanted developers with SQL skills... boosted it to No. 1."
           \end{itemize}
        \end{column}
        \begin{column}{0.38\textwidth}
            \begin{quote}
                "This ubiquity means that every software developer will have to interact with databases no matter the field, and SQL is the de facto standard..."
                \\ -- Prof. Andy Pavlo
            \end{quote}
        \end{column}
    \end{columns}
\end{frame}

%--- Slide 10 ---
\begin{frame}{Relational Languages}
    \begin{itemize}
        \item \textbf{DML:} Data Manipulation Language
        \item \textbf{DDL:} Data Definition Language
        \item \textbf{DCL:} Data Control Language
    \end{itemize}
    Also includes:
    \begin{itemize}
        \item View definition
        \item Integrity \& Referential Constraints
        \item Transactions
    \end{itemize}
    \begin{alertblock}{Important}
    SQL is based on \textbf{bags} (duplicates allowed) not \textbf{sets} (no duplicates).
    \end{alertblock}
\end{frame}

%--- Slide 11 ---
\begin{frame}{Today's Agenda}
    \begin{itemize}
        \item Aggregations + Group By
        \item String / Date / Time Operations
        \item Output Control + Redirection
        \item Window Functions
        \item Nested Queries
        \item Lateral Joins
        \item Common Table Expressions
    \end{itemize}
\end{frame}

%--- Slide 12 ---
\begin{frame}[fragile]{Example Database}
    
    \begin{columns}[t]
        \begin{column}{0.32\textwidth}
            \textbf{student} (sid, name, login, gpa)
            \begin{tabular}{|l|l|l|l|l|}
            \hline
            \textbf{sid} & \textbf{name} & \textbf{login} & \textbf{age} & \textbf{gpa} \\ \hline
            53666 & RZA & rza@cs & 55 & 4.0 \\ \hline
            53688 & Taylor & swift@cs & 27 & 3.9 \\ \hline
            53655 & Tupac & shakur@cs & 25 & 3.5 \\ \hline
            \end{tabular}
        \end{column}
        
        \begin{column}{0.32\textwidth}
            \textbf{enrolled} (sid, cid, grade)
            \begin{tabular}{|l|l|l|}
            \hline
            \textbf{sid} & \textbf{cid} & \textbf{grade} \\ \hline
            53666 & 15-445 & C \\ \hline
            53688 & 15-721 & A \\ \hline
            53688 & 15-826 & B \\ \hline
            53655 & 15-445 & B \\ \hline
            53666 & 15-721 & C \\ \hline
            \end{tabular}
        \end{column}
        
        \begin{column}{0.32\textwidth}
            \textbf{course} (cid, name)
            \begin{tabular}{|l|l|}
            \hline
            \textbf{cid} & \textbf{name} \\ \hline
            15-445 & Database Systems \\ \hline
            15-721 & Adv. DB Systems \\ \hline
            15-826 & Data Mining \\ \hline
            15-799 & Special Topics \\ \hline
            \end{tabular}
        \end{column}
    \end{columns}
\end{frame}

%--- Slide 13 ---
\begin{frame}{Aggregates}
    Functions that return a single value from a bag of tuples:
    \begin{itemize}
        \item \textbf{AVG(col)}: Return the average col value.
        \item \textbf{MIN(col)}: Return minimum col value.
        \item \textbf{MAX(col)}: Return maximum col value.
        \item \textbf{SUM(col)}: Return sum of values in col.
        \item \textbf{COUNT(col)}: Return \# of values for col.
    \end{itemize}
\end{frame}

%--- Slide 14-17 ---
\begin{frame}[fragile]{Aggregates}
    Aggregate functions can (almost) only be used in the SELECT output list.
    
    Get \# of students with a ``@cs'' login:
    
    \begin{lstlisting}
SELECT COUNT(login) AS cnt
FROM student WHERE login LIKE '%@cs'
    \end{lstlisting}
    
    \pause
    \begin{lstlisting}
SELECT COUNT(*) AS cnt
FROM student WHERE login LIKE '%@cs'
    \end{lstlisting}
    
    \pause
    \begin{lstlisting}
SELECT COUNT(1) AS cnt
FROM student WHERE login LIKE '%@cs'
    \end{lstlisting}
    
    \pause
    \begin{lstlisting}
SELECT COUNT(1+1+1) AS cnt
FROM student WHERE login LIKE '%@cs'
    \end{lstlisting}
\end{frame}

%--- Slide 18 ---
\begin{frame}[fragile]{Multiple Aggregates}
    Get the number of students and their average GPA that have a ``@cs'' login.
    
    \begin{lstlisting}
SELECT AVG(gpa), COUNT(sid)
FROM student WHERE login LIKE '%@cs'
    \end{lstlisting}
    
    \textbf{Result:}
    \begin{tabular}{ll}
    AVG(gpa) & COUNT(sid) \\
    3.8 & 3 
    \end{tabular}
\end{frame}

%--- Slide 19-20 ---
\begin{frame}[fragile]{Aggregates: Undefined Output}
    Output of other columns outside of an aggregate is undefined.
    
    Get the average GPA of students enrolled in each course.
    \begin{lstlisting}
SELECT AVG(s.gpa), e.cid
FROM enrolled AS e JOIN student AS s
  ON e.sid = s.sid
    \end{lstlisting}
    
    \begin{itemize}
        \item \texttt{AVG(s.gpa)} $\rightarrow$ 3.86
        \item \texttt{e.cid} $\rightarrow$ \textbf{??? (Undefined)}
    \end{itemize}
    
    \vspace{0.5cm}
    \textbf{Solution (Hack):}
    \begin{lstlisting}
SELECT AVG(s.gpa), ANY_VALUE(e.cid) ...
    \end{lstlisting}
\end{frame}

%--- Slide 21 ---
\begin{frame}[fragile]{Group By}
    Project tuples into subsets and calculate aggregates against each subset.
    
    \begin{columns}
        \begin{column}{0.5\textwidth}
        \begin{lstlisting}
SELECT AVG(s.gpa), e.cid
FROM enrolled AS e 
JOIN student AS s
  ON e.sid = s.sid
GROUP BY e.cid
        \end{lstlisting}
        \end{column}
        \begin{column}{0.4\textwidth}
            \scriptsize
            \textbf{Groups:}
            \begin{itemize}
                \item 15-721 (GPA: 2.25, 2.70) $\rightarrow$ Avg 2.46
                \item 15-826 (GPA: 2.75, 3.90, 3.50) $\rightarrow$ Avg 3.39
                \item 15-445 (GPA: 1.89) $\rightarrow$ Avg 1.89
            \end{itemize}
        \end{column}
    \end{columns}
\end{frame}

%--- Slide 22-24 ---
\begin{frame}[fragile]{Group By Requirements}
    Non-aggregated values in SELECT output clause \textbf{must} appear in GROUP BY clause.
    
    \textbf{Incorrect:}
    \begin{lstlisting}
SELECT AVG(s.gpa), e.cid, s.name
FROM enrolled AS e, student AS s
WHERE e.sid = s.sid
GROUP BY e.cid
    \end{lstlisting}
    
    \textbf{Correct:}
    \begin{lstlisting}
SELECT AVG(s.gpa), e.cid, s.name
...
GROUP BY e.cid, s.name
    \end{lstlisting}
\end{frame}

%--- Slide 25-27 ---
\begin{frame}[fragile]{Having}
    Filters results based on aggregation computation. (Like a WHERE clause for a GROUP BY).
    
    \begin{lstlisting}
SELECT AVG(s.gpa) AS avg_gpa, e.cid
FROM enrolled AS e, student AS s
WHERE e.sid = s.sid
GROUP BY e.cid
HAVING avg_gpa > 3.9;
    \end{lstlisting}
    
    \textit{Note: \texttt{WHERE} filters input tuples. \texttt{HAVING} filters output groups.}
\end{frame}

%--- Slide 28 ---
\begin{frame}{String Operations: Case Sensitivity}
    \begin{table}
        \begin{tabular}{l c c}
            \toprule
            \textbf{System} & \textbf{String Case} & \textbf{String Quotes} \\
            \midrule
            SQL-92 & Sensitive & Single Only \\
            Postgres & Sensitive & Single Only \\
            MySQL & Insensitive & Single/Double \\
            SQLite & Sensitive & Single/Double \\
            MSSQL & Sensitive & Single Only \\
            Oracle & Sensitive & Single Only \\
            \bottomrule
        \end{tabular}
    \end{table}
    \footnotesize
    Example: \texttt{WHERE UPPER(name) = UPPER('TuPaC')} vs \texttt{WHERE name = "TuPaC"}
\end{frame}

%--- Slide 29 ---
\begin{frame}[fragile]{String Operations: LIKE}
    \texttt{LIKE} is used for string matching.
    \begin{itemize}
        \item \textbf{'\%'} Matches any substring (including empty strings).
        \item \textbf{'\_'} Match any one character.
    \end{itemize}
    
    \begin{columns}
        \begin{column}{0.45\textwidth}
        \begin{lstlisting}
SELECT * FROM enrolled AS e
WHERE e.cid LIKE '15-%'
        \end{lstlisting}
        \end{column}
        \begin{column}{0.45\textwidth}
        \begin{lstlisting}
SELECT * FROM student AS s
WHERE s.login LIKE '%@c_'
        \end{lstlisting}
        \end{column}
    \end{columns}
\end{frame}

%--- Slide 30-31 ---
\begin{frame}[fragile]{String Operations: Functions}
    SQL-92 defines string functions, but DBMSs vary.
    
    \begin{lstlisting}
SELECT SUBSTRING(name,1,5) AS abbrv_name
FROM student WHERE sid = 53688
    \end{lstlisting}
    
    \textbf{Concatenation:}
    \begin{itemize}
        \item \textbf{SQL-92 / Postgres / Oracle:} \texttt{||} 
        \begin{lstlisting}
WHERE login = LOWER(name) || '@cs'
        \end{lstlisting}
        \item \textbf{MSSQL:} \texttt{+}
        \begin{lstlisting}
WHERE login = LOWER(name) + '@cs'
        \end{lstlisting}
        \item \textbf{MySQL:} \texttt{CONCAT}
        \begin{lstlisting}
WHERE login = CONCAT(LOWER(name), '@cs')
        \end{lstlisting}
    \end{itemize}
\end{frame}

%--- Slide 32 ---
\begin{frame}{Date/Time Operations}
    \begin{itemize}
        \item Operations to manipulate and modify DATE/TIME attributes.
        \item Can be used in both output and predicates.
        \item Support/syntax varies wildly between DBMSs.
    \end{itemize}
\end{frame}

%--- Slide 33-34 ---
\begin{frame}[fragile]{Output Redirection}
    Store query results in another table.
    
    \textbf{SQL-92 (Insert into existing):}
    \begin{lstlisting}
INSERT INTO CourseIds
(SELECT DISTINCT cid FROM enrolled);
    \end{lstlisting}
    
    \textbf{Create New Table (Varies):}
    \begin{itemize}
        \item \textbf{MySQL:} \texttt{CREATE TABLE CourseIds (SELECT ...)}
        \item \textbf{SQL-92 (Standard):} \texttt{SELECT ... INTO CourseIds FROM ...}
        \item \textbf{Postgres:} \texttt{SELECT ... INTO TEMPORARY CourseIds ...}
    \end{itemize}
\end{frame}

%--- Slide 35-37 ---
\begin{frame}[fragile]{Output Control: ORDER BY}
    \texttt{ORDER BY <column*> [ASC|DESC]}
    
    \begin{lstlisting}
SELECT sid, grade FROM enrolled
WHERE cid = '15-721'
ORDER BY grade
    \end{lstlisting}
    
    Can also order by multiple columns or column index:
    \begin{lstlisting}
ORDER BY grade DESC, sid ASC
    \end{lstlisting}
    
    \begin{lstlisting}
ORDER BY 2  -- Sort by second column (grade)
    \end{lstlisting}
\end{frame}

%--- Slide 38-40 ---
\begin{frame}[fragile]{Output Control: LIMIT / FETCH}
    Limit the \# of tuples returned.
    
    \textbf{Standard Syntax:}
    \begin{lstlisting}
SELECT sid, name FROM student
WHERE login LIKE '%@cs'
FETCH FIRST 10 ROWS ONLY;
    \end{lstlisting}
    
    \textbf{With Offset (Range):}
    \begin{lstlisting}
ORDER BY gpa
OFFSET 10 ROWS
FETCH FIRST 10 ROWS WITH TIES;
    \end{lstlisting}
\end{frame}

%--- Slide 41-43 ---
\begin{frame}{Window Functions}
    Performs a calculation across a set of tuples that are related to the current tuple, without collapsing them into a single output tuple.
    
    \begin{itemize}
        \item Supports running totals, ranks, moving averages.
        \item Syntax: \texttt{SELECT FUNC-NAME(...) OVER (...) FROM tableName}
    \end{itemize}
    
    \textbf{Functions:}
    \begin{itemize}
        \item Aggregation functions (AVG, MIN, MAX...)
        \item \textbf{ROW\_NUMBER()}: \# of the current row
        \item \textbf{RANK()}: Order position of the current row
    \end{itemize}
\end{frame}

%--- Slide 44 ---
\begin{frame}[fragile]{Window Functions: Example}
    \begin{lstlisting}
SELECT *, ROW_NUMBER() OVER () AS row_num
FROM enrolled
    \end{lstlisting}
    
    \begin{tabular}{|l|l|l|l|}
    \hline
    sid & cid & grade & row\_num \\ \hline
    53666 & 15-445 & C & 1 \\ \hline
    53688 & 15-721 & A & 2 \\ \hline
    53688 & 15-826 & B & 3 \\ \hline
    ... & ... & ... & ... \\ \hline
    \end{tabular}
\end{frame}

%--- Slide 45-47 ---
\begin{frame}[fragile]{Window Functions: PARTITION BY}
    The \texttt{OVER} keyword specifies how to group together tuples when computing the window function.
    
    \begin{lstlisting}
SELECT cid, sid,
  ROW_NUMBER() OVER (PARTITION BY cid)
FROM enrolled
ORDER BY cid
    \end{lstlisting}
    
    \textbf{Result:} Resets counter for each \texttt{cid}.
    \begin{tabular}{llc}
    \textbf{cid} & \textbf{sid} & \textbf{row\_number} \\
    15-445 & 53666 & 1 \\
    15-445 & 53655 & 2 \\
    15-721 & 53688 & 1 \\
    \end{tabular}
\end{frame}

%--- Slide 48 ---
\begin{frame}[fragile]{Window Functions: ORDER BY}
    You can also include an \texttt{ORDER BY} in the window grouping to sort entries in each group.
    
    \begin{lstlisting}
SELECT *,
  ROW_NUMBER() OVER (ORDER BY cid)
FROM enrolled
ORDER BY cid
    \end{lstlisting}
\end{frame}

%--- Slide 49-50 ---
\begin{frame}[fragile]{Window Functions: Filtering Ranks}
    Find the student with the second highest grade for each course.
    
    \begin{lstlisting}
SELECT * FROM (
  SELECT *, RANK() OVER (PARTITION BY cid
                         ORDER BY grade ASC) AS rank
  FROM enrolled
) AS ranking
WHERE ranking.rank = 2
    \end{lstlisting}
    \begin{enumerate}
        \item Group tuples by cid
        \item Then sort by grade
    \end{enumerate}
\end{frame}

%--- Slide 51-53 ---
\begin{frame}[fragile]{Nested Queries}
    Invoke a query inside of another query to compose more complex computations.
    
    \begin{lstlisting}
SELECT name FROM student WHERE
sid IN (SELECT sid FROM enrolled)
    \end{lstlisting}
    
    \begin{itemize}
        \item Outer Query vs Inner Query.
        \item Can appear in SELECT output, FROM clause, or WHERE clause.
    \end{itemize}
\end{frame}

%--- Slide 54-57 ---
\begin{frame}[fragile]{Nested Queries Example}
    Get the names of students in '15-445'.
    
    \begin{lstlisting}
SELECT name FROM student
WHERE sid IN (
    SELECT sid FROM enrolled
    WHERE cid = '15-445'
)
    \end{lstlisting}
\end{frame}

%--- Slide 58 ---
\begin{frame}{Nested Queries: Operators}
    \begin{itemize}
        \item \textbf{ALL}: Must satisfy expression for \textit{all} rows in sub-query.
        \item \textbf{ANY}: Must satisfy expression for \textit{at least one} row in sub-query.
        \item \textbf{IN}: Equivalent to \texttt{=ANY()}.
        \item \textbf{EXISTS}: At least one row is returned (no attribute comparison needed).
    \end{itemize}
\end{frame}

%--- Slide 59 ---
\begin{frame}[fragile]{Nested Queries: ANY Example}
    Get the names of students in '15-445':
    
    \begin{lstlisting}
SELECT name FROM student
WHERE sid = ANY(
    SELECT sid FROM enrolled
    WHERE cid = '15-445'
)
    \end{lstlisting}
\end{frame}

%--- Slide 60-63 ---
\begin{frame}[fragile]{Nested Queries: Highest ID}
    Find student record with the highest id that is enrolled in at least one course.
    
    \begin{lstlisting}
SELECT sid, name FROM student
WHERE sid >= ALL(
    SELECT sid FROM enrolled
)
    \end{lstlisting}
    
    \textit{(Alternatively using MAX)}
    \begin{lstlisting}
SELECT sid, name FROM student
WHERE sid IN (
    SELECT MAX(sid) FROM enrolled
)
    \end{lstlisting}
\end{frame}

%--- Slide 64-65 ---
\begin{frame}[fragile]{Nested Queries: NOT EXISTS}
    Find all courses that have no students enrolled in it.
    
    \begin{lstlisting}
SELECT * FROM course
WHERE NOT EXISTS(
    SELECT * FROM enrolled
    WHERE course.cid = enrolled.cid
)
    \end{lstlisting}
\end{frame}

%--- Slide 66 ---
\begin{frame}[fragile]{Lateral Joins}
    The \textbf{LATERAL} operator allows a nested query to reference attributes in other nested queries that precede it.
    
    \begin{itemize}
        \item Think of it like a \texttt{for} loop.
    \end{itemize}
    
    \begin{lstlisting}
SELECT * FROM
  (SELECT 1 AS x) AS t1,
  LATERAL (SELECT t1.x+1 AS y) AS t2;
    \end{lstlisting}
    \textbf{Result:} x=1, y=2
\end{frame}

%--- Slide 67-69 ---
\begin{frame}[fragile]{Lateral Join Example}
    Calculate the number of students enrolled in each course and the average GPA.
    
    \begin{lstlisting}
SELECT * FROM course AS c,
  LATERAL (SELECT COUNT(*) AS cnt FROM enrolled
           WHERE enrolled.cid = c.cid) AS t1,
  LATERAL (SELECT AVG(gpa) AS avg FROM student AS s
           JOIN enrolled AS e ON s.sid = e.sid
           WHERE e.cid = c.cid) AS t2;
    \end{lstlisting}
\end{frame}

%--- Slide 70-71 ---
\begin{frame}[fragile]{Common Table Expressions (CTE)}
    Specify a temporary result set that can then be referenced by another part of that query.
    \begin{itemize}
        \item Alternative to nested queries and views.
    \end{itemize}
    
    \begin{lstlisting}
WITH cteName AS (
    SELECT 1
)
SELECT * FROM cteName
    \end{lstlisting}
\end{frame}

%--- Slide 72-74 ---
\begin{frame}[fragile]{CTE Binding}
    You can bind/alias output columns to names before the AS keyword.
    
    \begin{lstlisting}
WITH cteName (col1, col2) AS (
    SELECT 1, 2
)
SELECT col1 + col2 FROM cteName
    \end{lstlisting}
    
    \textbf{Example (Max ID):}
    \begin{lstlisting}
WITH cteSource (maxId) AS (
    SELECT MAX(sid) FROM enrolled
)
SELECT name FROM student, cteSource
WHERE student.sid = cteSource.maxId
    \end{lstlisting}
\end{frame}

%--- Slide 75 ---
\begin{frame}{Other Things To Note}
    \begin{itemize}
        \item Identifiers (table and column names) are case-insensitive.
        \item Makes it harder for applications that care about case (e.g., CamelCase).
        \item One often sees quotes around names: \texttt{"ArtistList.firstName"}
        \item \textit{ISO/IEC 9075-2:2023 Standard Document costs CHF 216.}
    \end{itemize}
\end{frame}

%--- Slide 76 ---
\begin{frame}{Conclusion}
    \begin{columns}
        \begin{column}{0.5\textwidth}
            \begin{itemize}
                \item SQL is a hot language.
                \item Lots of NL2SQL tools, but writing SQL is not going away.
                \item You should (almost) always strive to compute your answer as a single SQL statement.
            \end{itemize}
        \end{column}
        \begin{column}{0.5\textwidth}
            \textbf{Top Programming Languages 2023 (IEEE Spectrum)}
            \begin{table}
                \begin{tabular}{lr}
                \textbf{Language} & \textbf{Rank} \\ \hline
                \rowcolor{yellow} SQL & 1.00 \\
                Python & 0.99 \\
                Java & 0.76 \\
                JavaScript & 0.55 \\
                C++ & 0.45 \\
                \end{tabular}
            \end{table}
        \end{column}
    \end{columns}
\end{frame}

%--- Slide 77 ---
\begin{frame}{Homework \#1}
    \textbf{Write SQL queries to perform basic data analysis.}
    \begin{itemize}
        \item Write the queries locally using SQLite + DuckDB.
        \item Submit them to Gradescope.
        \item You can submit multiple times and use your best score.
    \end{itemize}
    \vspace{0.5cm}
    \textbf{Due: Sunday Sept 8th @ 11:59pm}
\end{frame}

%--- Slide 78 ---
\begin{frame}{Next Class}
    \centering
    \Large
    We will begin our journey to understanding the internals of database systems starting with \textbf{Storage}!
\end{frame}

\end{document}